\documentclass[12pt]{article}
\usepackage[T1]{fontenc}
\usepackage[brazilian]{babel}
\usepackage{csquotes}
\usepackage[
  perfil=artigo,
  artigo-colunas=duas,
  citacoes=autor-data
]{abntex3}

\addbibresource{referencias-exemplo.bib}

\ABNTEXsetup{
  titulo={Preservação de artigos científicos digitais},
  subtitulo={um fluxo comunitário reproduzível}
}
\ABNTEXartigoSetup{
  titulo-outro-idioma={Preservation of digital scientific articles},
  resumo={Este artigo descreve um fluxo reproduzível para preservar documentos científicos digitais.},
  palavras-chave={preservação digital; normalização; LaTeX},
  resumo-outro-idioma={This article describes a reproducible workflow for preserving digital scientific documents.},
  palavras-chave-outro-idioma={digital preservation; standardization; LaTeX},
  data-submissao={10 de janeiro de 2026},
  data-aprovacao={20 de março de 2026},
  identificacao-disponibilidade={Disponível em: \texttt{https://example.org/artigo}}
}
\ABNTEXartigoautor
  {Ana Silva}
  {Doutora em Ciência da Informação}
  {Universidade Exemplo}
  {ana@example.org}
\ABNTEXartigoautor
  {Bruno Souza}
  {Mestre em Preservação Digital}
  {Instituto Exemplo}
  {bruno@example.org}

\begin{document}
\ABNTEXartigopretextual
\ABNTEXartigotextual

\ABNTEXartigointroducao{%
  A preservação digital exige procedimentos verificáveis. A apresentação
  documental também se apoia em normas técnicas \parencite{associacao2025}.
}
\ABNTEXartigodesenvolvimento[Procedimentos e resultados]{%
  O fluxo mantém fontes, referências e artefatos em repositório versionado,
  executa testes automatizados e valida a distribuição isoladamente.
}
\ABNTEXartigoconsideracoesfinais{%
  A automação reduz falhas de empacotamento e favorece a reprodução.
}

\ABNTEXartigopostextual
\ABNTEXbibliografia
\ABNTEXartigoagradecimentos{À comunidade que revisou o fluxo técnico.}
\end{document}
